\section{From Dedicated Storage Devices to Archival Manufacture}\label{sec:intro}

Digital knowledge currently survives through a chain of dependencies that are rarely examined together: functioning organizations that host it, networks that deliver it, accounts that authorize access to it, electricity that powers the machines that hold it, and specialized storage equipment whose formats are understood by a small number of manufacturers. Each link is individually reliable, and the chain as a whole has proved durable for a few decades, but it is a chain, and a failure at any link removes access for everyone downstream. The proposal examined here concerns a different mode of survival. Selected public archives would travel not only through dedicated storage devices but also through manufacturing and construction, so that reference collections might be embodied in tiles, technical manuals embedded in replacement parts, and dated encyclopedic snapshots distributed with household objects. The archive would then persist for the same reason that a pattern in a brick persists, namely that the object is useful for other reasons and is therefore made, shipped, installed, and kept.

The motivating object is an ordinary brick carrying two inscriptions. The first is an obvious and sparse introductory inscription that states what the object contains and how to read it. The second is a dense optical archive written into a face or into the volume of the unit during manufacture. This single example poses the entire problem at once. The archive must survive manufacture, which subjects it to heat, pressure, shrinkage, and handling. It must remain discoverable by someone who does not know it is there, which requires that the introductory layer be recognizable. It must permit alignment and calibration of a reader to the particular specimen, since no two bricks are identical. It must identify its own format. And it must yield verified files through a reader that an independent party can reconstruct from published specifications, because an archive that can be read only with a proprietary scanner, a cloud account, or an undocumented decoding model has replaced one dependency with another.

The value of such a scheme has two distinct sources that the essay keeps apart throughout. The first is density, meaning the number of recoverable bits per unit area or volume of material. The second is distribution, meaning the number of independent physical copies and the diversity of the objects and places that hold them. A modest archive replicated across millions of ordinary objects could be socially consequential even if its density never exceeded that of specialized archival media, because the consequence derives from the cost of replication and the dispersal of copies rather than from the packing of bits. Conversely, an extraordinarily dense medium confined to a few vaults provides the vulnerability of centralized storage regardless of how much each vault holds. Density determines how large an archive an object of a given size can carry, and distribution determines how likely it is that some copy survives, and the program must establish both.

The technical argument is organized as follows. Section~\ref{sec:foundations} situates the proposal among existing work on laser-written glass, volume holography, and multiplexed optical encoding, and identifies what a new contribution would have to be. Section~\ref{sec:defs} fixes the vocabulary that separates addresses from payloads and fabrication controls from observable states. Section~\ref{sec:model} introduces the write--store--read model on which every later claim rests. Sections~\ref{sec:capacity} through~\ref{sec:fixed} develop the mathematical core: the combinatorial capacity of a product alphabet, geometric scaling to areas and volumes, the limiting role of observational equivalence, robust capacity as a packing problem, overlapping storage, linear multiplexing with crosstalk, and fixed-resource scaling in a Gaussian channel. Sections~\ref{sec:substrates} through~\ref{sec:ecc} address substrates, reader circuits, calibration and self-description, and error correction. Section~\ref{sec:stego} treats concealment and confidentiality, which for an archive intended to be publicly recoverable are in tension with the principal design goal and therefore require careful accounting. Sections~\ref{sec:experiments} through~\ref{sec:distribution} specify the experimental sequence, the controls, and the system-level analysis of replication. Section~\ref{sec:popular} returns to the broader vision in a less formal register, and Section~\ref{sec:program} lays out the resulting publication program. Appendices collect worked budgets, a container format, and the notation.

\section{Existing Foundations and the Proposed Contribution}\label{sec:foundations}

The proposal stands on substantial prior engineering, and the purpose of this section is to identify that foundation precisely enough that the contribution of the present program can be stated without overclaiming. Laser writing in glass is the nearest precedent. Early demonstrations of five-dimensional recording in glass combined three spatial coordinates with two optical properties of laser-induced nanostructures, namely the orientation and the strength of the retardance that the structure induces in transmitted light, and showed that polarization-controlled material modifications could carry digital information \cite{zhang2013}. Subsequent work demonstrated multilayer recording with error-free recovery over many layers in silica \cite{wang2022}. The most developed recent result, reported by the Project Silica team, integrates inscription, optical readout, decoding, error correction, and retention testing, and reports usable capacities of several terabytes in a single glass specimen of standard plate dimensions \cite{silica2026}. The accompanying systems literature considers how such a medium becomes part of an archival architecture rather than remaining a laboratory demonstration \cite{anderson2023}. These results establish that volumetric optical archiving in glass is a real engineering discipline with measured capacities. They do not establish that comparable performance transfers to stone, ceramic, cloth, or the heterogeneous materials of the built environment, and the present program treats that transfer as an open experimental question.

A second body of work concerns holographic data storage, in which data are encoded as interference patterns and multiple pages share a volume through angular, wavelength, or phase multiplexing. The foundational reference collects the theory of recording media, multiplexing, optical architectures, and system optimization \cite{coufal2000}, and a recent survey examines the density, energy efficiency, and switching difficulties that have repeatedly prevented holographic storage from displacing competing technologies \cite{cheriere2025}. That survey is an essential counterweight to capacity arguments based solely on nominal degrees of freedom, because it records the dynamic-range and interference limits that appear when many holograms share a material.

A third body of work concerns fabricated metasurfaces, in which structure at or below the wavelength scale controls the phase, amplitude, and polarization of reflected or transmitted light. Metasurface holography has demonstrated orbital-angular-momentum multiplexing in momentum space with complex-amplitude control \cite{ren2020}, and structural multiplexing with diffractive neural networks in three-dimensionally printed elements \cite{fan2024}. Those demonstrations concern reconstructed images, and image fidelity should not be equated with verified digital-file capacity, since the error criterion for an archive is exact recovery of every bit after error correction rather than perceptual similarity.

Optical research already investigates multiplexing through wavelength, polarization, incident angle, and orbital angular momentum, so the general observation that light supplies many usable degrees of freedom is not a contribution. The contribution of the present program, if it has one, must be located in one of four places: a specific encoding architecture whose capacity is established under an explicit error criterion; a fabrication method that writes the encoding during ordinary manufacture rather than in a dedicated inscription step; a reader that can be reconstructed independently from published specifications; or an archival system that makes replication across everyday objects operationally meaningful. The essay develops all four at the level of specification and bounds, and the experimental sections state which claims each experiment would support or refute.

The mathematical apparatus is deliberately elementary, since the central difficulty lies in stating the counting problem correctly rather than in advanced technique. Counting principles, functions and relations, equivalence classes, and the pigeonhole principle, as presented in standard discrete mathematics texts \cite{epp,kolman}, suffice for the first half of the development. Shannon's framework for reliable communication through noisy channels \cite{shannon1948,cover2006} supplies the second half. Digital control and circuit analysis are treated using the standard treatments of logic design \cite{mano} and linear circuits \cite{boylestad}, supplemented in Section~\ref{sec:reader} by specialized references that a final bibliography would add for photodetection and optical design.

\section{Definitions: What Is Being Counted}\label{sec:defs}

Most errors in popular discussions of multidimensional storage arise from counting the wrong thing, and the definitions below are chosen to make the intended count explicit. A \emph{manufacturing control} is an adjustable input to the fabrication process, such as pulse duration, pulse energy, inscription orientation, feature depth, spatial frequency, layer position, or composition. A \emph{material state} is the physical configuration that results from fabrication, which depends on the controls and also on the substrate and on process variation. A \emph{readout observable} is a measurable response of a material state under a specified reader configuration, such as transmitted intensity, retardance, scattering angle, or phase. A \emph{recoverable symbol} is a member of a codebook whose elements are distinguishable under stated operating conditions with a stated error probability. These four notions form a chain from what the manufacturer adjusts to what the reader can decide, and a count taken at the first link says nothing about a count taken at the last.

Three distinctions follow, each of which prevents a specific accounting error. The first is between controls and states: several different settings of the controls can produce the same material state, as when two pulse widths leave the same mark, so the number of control settings overstates the number of states. The second is between states and observations: different material states can produce indistinguishable observations, so the number of states overstates the number of recoverable symbols. The third is between an address and a payload: an enormous space of possible inscriptions is an address space, and it becomes a storage capacity only to the extent that each address can be written and read reliably and carries information that the reader did not already possess. A physical storage \emph{site} is a region that is written and read as a unit, and it must be distinguished from an optical \emph{mode}, which is one of several independently addressable components of the response of a single site. An overlapping archive, in which several collections share a region, must in turn be distinguished from an alternative interpretation of one archive, in which the same observations are decoded in different ways to yield different files, since only the former adds payload.

\begin{definition}[Capacity of a configuration space]\label{def:capacity}
Let $\mathcal U$ be a finite set of fabrication settings, $\mathcal Y$ the set of readout observations, and $\omega:\mathcal U\to\mathcal Y$ the nominal readout map. A \emph{zero-error codebook} is a subset $\mathcal K\subseteq\mathcal U$ on which $\omega$ is injective, and the \emph{zero-error capacity} of one site is $\lfloor\log_2 |\omega(\mathcal U)|\rfloor$ bits, which equals the logarithm of the number of distinct observations that the reader can produce.
\end{definition}

The definition has a consequence that will be used repeatedly: capacity is a property of the image of the readout map and not of its domain. The number of fabrication settings, which is the quantity that grows without limit as parameters are added, enters the capacity only through the number of distinct observations those settings produce.

\section{The Complete Write--Store--Read Model}\label{sec:model}

Every capacity claim in the essay refers to the following chain of operations, which makes the roles of substrate variation, ageing, and reader configuration explicit:
\[
X\xrightarrow{\;E\;}U\xrightarrow{\;F_s\;}Z_0\xrightarrow{\;D_t\;}Z_t\xrightarrow{\;R_v\;}Y\xrightarrow{\;G\;}\widehat X .
\]
Here $X$ is the source message, $E$ the encoder, which includes compression, error-correcting coding, and the mapping from coded symbols to fabrication settings, $U$ the resulting fabrication instructions, $F_s$ the substrate-dependent fabrication process indexed by a vector $s$ of substrate and process conditions, $Z_0$ the as-fabricated material state, $D_t$ the degradation process over time $t$, $Z_t$ the aged state, $R_v$ the reader operating under configuration $v$, $Y$ the raw observation, $G$ the decoder, and $\widehat X$ the recovered message. The maps $F_s$, $D_t$, and $R_v$ are random in general, and the model should be read as a composition of channels rather than of deterministic functions.

The decomposition matters because failures belong to different links and have different remedies. A decoder may work on an ideal simulation and fail after printing shrinkage, which is a failure of $F_s$. A physical inscription may remain stable for centuries while its reader becomes unavailable, which is a failure of $R_v$ and of the documentation that would permit its reconstruction. An inscription that is fully readable on the day of manufacture may be unreadable after abrasion or weathering, which is a failure of $D_t$. A program that measured only the first link would overstate the usefulness of an archive, and all three kinds of failure belong within the research program.

The primary performance condition is
\begin{equation}\label{eq:perr}
\Pro\bigl(\widehat X\neq X\bigr)\le\varepsilon,
\end{equation}
evaluated under a specified distribution of manufacturing conditions $s$, environmental exposures $t$, and reader tolerances $v$. Every capacity claim in the essay is accompanied by a statement of $\varepsilon$ and of that distribution, because a capacity without an error criterion is not a well-defined quantity. For an archive of verified files the natural criterion is exact recovery of each block after error correction, with a small $\varepsilon$ per block and an explicit policy for blocks that cannot be recovered.
